\section*{الفصل \ref{ch:b1:lewis-resonance-vsepr} --- بنى لويس والرنين ونموذج VSEPR}

\begin{solution}{exo:b1:lewis-resonance-vsepr:1}
\ce{H2O}: رابطتان \ce{O-H}، \omterm{def:b1:lewis-resonance-vsepr:lewis}{وزوجان حران} على O. \ce{NH3}: ثلاث
روابط \ce{N-H}، \omterm{def:b1:lewis-resonance-vsepr:lewis}{وزوج حر} واحد على N. \ce{CO2}: \ce{O=C=O}، \omterm{def:b1:lewis-resonance-vsepr:lewis}{وزوجان حران}
على كل O. \ce{HCN}: \ce{H-C#N}، \omterm{def:b1:lewis-resonance-vsepr:lewis}{وزوج حر} واحد على N. \ce{CH2O}: الكربون
مرتبط بذرتي H وبرابطة ثنائية مع O، \omterm{def:b1:lewis-resonance-vsepr:lewis}{وزوجان حران} على O. كل \omterm{def:b1:lewis-resonance-vsepr:formal-charge}{الشحنات
الشكلية} معدومة في هذه الخمسة. \ce{NH4+}: أربع روابط \ce{N-H}، ولا \omterm{def:b1:lewis-resonance-vsepr:lewis}{زوج
حر}؛ النيتروجين $5 - 0 - 4 = +1$، وهي شحنة الأيون.
\end{solution}

\begin{solution}{exo:b1:lewis-resonance-vsepr:2}
\ce{H2S}: \ce{AX2E2}، منحنٍ، دون $109.5^\circ$ (المقيسة $92^\circ$).
\ce{PCl3}: \ce{AX3E}، هرمي ثلاثي، نحو $100^\circ$.
\ce{BF3}: \ce{AX3}، مثلثي مستوٍ، $120^\circ$. \ce{SiH4}: \ce{AX4}،
رباعي الأوجه، $109.5^\circ$. \ce{CO3^{2-}}: \ce{AX3} (الرابطة الثنائية
تُعدّ مرة واحدة)، مثلثي مستوٍ، $120^\circ$.
% ledger: geo:H2S.aHSH, geo:PCl3.aClPCl
\end{solution}

\begin{solution}{exo:b1:lewis-resonance-vsepr:3}
غير قطبية: \ce{CCl4} (\ce{AX4} رباعي الأوجه)، \ce{BF3} (\ce{AX3})،
\ce{CO2} (\ce{AX2} خطي): تتلاشى ثنائيات أقطاب الروابط. قطبية: \ce{CH2Cl2}
(نوعان مختلفان من الروابط على رباعي أوجه)، \ce{NH3} (هرمي)،
\ce{SO2} (منحنٍ).
\end{solution}

\begin{solution}{exo:b1:lewis-resonance-vsepr:4}
\ce{H+} هو \omterm{def:b1:lewis-resonance-vsepr:lewis-acid}{حمض لويس} و\ce{H2O} القاعدة: الرابطة \ce{O-H} الجديدة في
\ce{H3O+} تناسقية عند تكوّنها (ثم تصير روابط \ce{O-H} الثلاث
متماثلة). و\ce{Cu^{2+}} هو الحمض وكل \ce{NH3} قاعدة: الروابط \ce{Cu-N}
الأربع تناسقية.
\end{solution}

\begin{solution}{exo:b1:lewis-resonance-vsepr:5}
\babelsublr{\chemfig{CH_3-C(=[:60]O)-[:-60]\chemabove{O}{\scriptstyle\ominus}}
$\leftrightarrow$
\chemfig{CH_3-C(-[:60]\chemabove{O}{\scriptstyle\ominus})=[:-60]O}}.
البنيتان متكافئتان: لكل من الرابطتين \ce{C-O} الرتبة 1.5
والطول نفسه، بين الرابطة الثنائية والرابطة الأحادية. أما في حمض الميثانويك
فذرتا الأكسجين مختلفتان (إحداهما تحمل H)، فالبنيتان غير
متكافئتين، وتحتفظ الرابطتان بطولين متمايزين، 120.2 (\ce{C=O})
و\qty{134.3}{pm} (\ce{C-OH}).
\end{solution}

\begin{solution}{exo:b1:lewis-resonance-vsepr:6}
$N = 6 + 4 \times 6 + 2 = 32$ إلكترونًا. بأربع روابط أحادية وثلاثة
\omterm{def:b1:lewis-resonance-vsepr:lewis}{أزواج حرة} على كل أكسجين: الكبريت $6 - 0 - 4 = +2$، وكل أكسجين
$6 - 6 - 1 = -1$؛ المجموع $+2 - 4 = -2$. الشكل \ce{AX4}، رباعي الأوجه. وبرابطتين
\ce{S=O}، يعطي اختيار ذرتي الأكسجين المرتبطتين برابطة ثنائية من بين
الأربع $\binom{4}{2} = 6$ بنى.
\end{solution}

\begin{solution}{exo:b1:lewis-resonance-vsepr:7}
\ce{SF4}: \ce{AX4E}، أرجوحي. \ce{ClF3}: \ce{AX3E2}، على شكل حرف T.
\ce{XeF2}: \ce{AX2E3}، خطي. \ce{XeF4}: \ce{AX4E2}، مربع مستوٍ. في
الهرم الثلاثي المزدوج للموضع المحوري ثلاثة جيران عند $90^\circ$،
وللموضع الاستوائي جاران فقط: فتذهب \omterm{def:b1:lewis-resonance-vsepr:lewis}{الأزواج الحرة} الضخمة حيث تلقى
أقل عدد من التنافرات عند $90^\circ$.
\end{solution}

\begin{solution}{exo:b1:lewis-resonance-vsepr:8}
الكبريت والفوسفور أكبر وأقل كهرسلبية من الأكسجين
والنيتروجين: فتقع \omterm{def:b1:lewis-resonance-vsepr:lewis}{الأزواج الرابطة} أبعد عن الذرة المركزية وتنجذب
نحو الهيدروجين، فيقل تنافرها، وتضغط \omterm{def:b1:lewis-resonance-vsepr:lewis}{الأزواج
الحرة} الروابط نحو $90^\circ$.
\end{solution}

\begin{solution}{exo:b1:lewis-resonance-vsepr:9}
\ce{CH3-C#N}: كربون \ce{CH3} من النوع \babelsublr{sp$^3$}، وكربون النتريل sp،
والنيتروجين sp؛ 5 روابط $\sigma$ (ثلاث \ce{C-H}، و\ce{C-C}، وواحدة في
\ce{C#N}) ورابطتا $\pi$. \ce{CH2=CH-CH3}: كربونا الألكين
\babelsublr{sp$^2$}، وكربون الميثيل \babelsublr{sp$^3$}؛ 8 روابط $\sigma$ (ست \ce{C-H}، واثنتان
\ce{C-C}) ورابطة $\pi$ واحدة.
\end{solution}

\begin{solution}{exo:b1:lewis-resonance-vsepr:10}
الجزيئان كلاهما هرميان \omterm{def:b1:lewis-resonance-vsepr:lewis}{بزوج حر} على النيتروجين. ويعطي \omterm{def:b1:lewis-resonance-vsepr:lewis}{الزوج الحر}
إسهامًا متجهًا (من السالب إلى الموجب) من \omterm{def:b1:lewis-resonance-vsepr:lewis}{الزوج الحر}
نحو النواة. في \ce{NH3} يكون النيتروجين الطرف السالب لكل
رابطة: فتتجه ثنائيات أقطاب الروابط من N نحو الهيدروجينات، في
اتجاه إسهام \omterm{def:b1:lewis-resonance-vsepr:lewis}{الزوج الحر} نفسه، فتتجمع. وفي \ce{NF3}
يكون الفلور الطرف السالب: فتتجه ثنائيات أقطاب الروابط من الفلورات
نحو N، عكس إسهام \omterm{def:b1:lewis-resonance-vsepr:lewis}{الزوج الحر}، فيكاد الاثنان يتلاشيان.
\end{solution}

\begin{solution}{exo:b1:lewis-resonance-vsepr:11}
$\mu_b = \mu/(2\cos(\theta/2)) = 1.63/(2\cos 59.75^\circ) =
\qty{1.62}{\debye} = \qty{5.40e-30}{C.m}$. ومنه $\delta = \mu_b/(e\,d) =
5.40 \times 10^{-30}/(1.602 \times 10^{-19} \times 143.2 \times 10^{-12})
= 0.24$.
\end{solution}

\begin{solution}{exo:b1:lewis-resonance-vsepr:12}
\ce{N=N=O} \omterm{def:b1:lewis-resonance-vsepr:formal-charge}{بالشحنات الشكلية} $-1$، $+1$، $0$؛ و\ce{N#N-O} بالشحنات $0$،
$+1$، $-1$. الرابطة \ce{N-N} (\qty{112.8}{pm}) قريبة من الرابطة
الثلاثية في \ce{N2} (\qty{109.8}{pm}): فالبنية \ce{N#N-O} أثقل وزنًا،
وهي تضع الشحنة السالبة على الأكسجين، الذرة الأكثر
كهرسلبية. غير أن الرابطة \ce{N-O} أقصر من رابطة أحادية: فالبنية
الأخرى تسهم أيضًا.
\end{solution}

\begin{solution}{pb:b1:lewis-resonance-vsepr:1}
\textbf{1.} \ce{N2} 10، \ce{NO} 11، \ce{NO2} 17، \ce{N2O} 16، \ce{O3} 18،
\ce{HNO3} 24.
\textbf{2.} \ce{N#N} مع \omterm{def:b1:lewis-resonance-vsepr:lewis}{زوج حر} واحد على كل نيتروجين.
\textbf{3.} \ce{N=O} مع \omterm{def:b1:lewis-resonance-vsepr:lewis}{زوجين حرين} على الأكسجين، \omterm{def:b1:lewis-resonance-vsepr:lewis}{وزوج حر} واحد
\omterm{def:b1:quantum-numbers:unpaired}{وإلكترون منفرد} على النيتروجين. ومع 11 إلكترونًا، وهو عدد فردي، يكون لإحدى الذرتين
دائمًا عدد فردي: للنيتروجين 7 إلكترونات حوله.
\textbf{4.} \ce{N=N=O}: $-1$، $+1$، $0$؛ \ce{N#N-O}: $0$، $+1$، $-1$.
\textbf{5.} \ce{O=O-O}: الأكسجين المركزي $+1$ (\omterm{def:b1:lewis-resonance-vsepr:lewis}{زوج حر} واحد، وثلاث روابط)،
والأكسجين الطرفي المرتبط برابطة أحادية $-1$، والآخر 0.
\textbf{6.} كما في \cref{ex:b1:lewis-resonance-vsepr:hno3}: النيتروجين $+1$،
والأكسجين الحامل لرابطة أحادية دون هيدروجين $-1$، والبقية 0.
\textbf{7.} \ce{NO} و\ce{NO2} (11 و17 إلكترونًا).
\textbf{8.} بنيتان متكافئتان، والرابطة الثنائية في هذا الجانب أو ذاك:
الرتبة 1.5 لكل رابطة.
\textbf{9.} نعم: تقع \qty{127.8}{pm} بين الرابطة الثنائية
(\qty{120.8}{pm}) والرابطة الأحادية (\qty{147.5}{pm}).
\textbf{10.} الرابطة الثنائية على كل من ذرات الأكسجين الثلاث بالتناوب؛ فكل
رابطة \ce{N-O} ثنائية في بنية واحدة من ثلاث: الرتبة $4/3$.
\textbf{11.} الرابطة الطويلة، \qty{140.6}{pm}، هي \ce{N-OH}، رابطة أحادية.
وتتقاسم ذرتا الأكسجين الخاليتان من الهيدروجين رابطة $\pi$ واحدة عبر \omterm{def:b1:lewis-resonance-vsepr:resonance}{بنيتي
رنين} متكافئتين: رابطتان من الرتبة 1.5، 121.1
و\qty{119.9}{pm}.
\textbf{12.} يحمل النيتروجين إلكترونًا منفردًا لا \omterm{def:b1:lewis-resonance-vsepr:lewis}{زوجًا حرًا}:
فهو يدفع \omterm{def:b1:lewis-resonance-vsepr:lewis}{الأزواج الرابطة} أقل مما يدفعها زوج، فتنفرج الروابط
إلى ما وراء $120^\circ$.
\textbf{13.} \ce{N#N-O}، والشحنة السالبة على الأكسجين.
\textbf{14.} \ce{O3}: \ce{AX2E}، منحنٍ. \ce{N2O}: \ce{AX2} (N مركزي)،
خطي. \ce{NO3-}: \ce{AX3}، مثلثي مستوٍ. نيتروجين \ce{HNO3}:
\ce{AX3}، مثلثي مستوٍ.
\textbf{15.} \omterm{def:b1:lewis-resonance-vsepr:lewis}{الزوج الحر} للأكسجين المركزي يدفع \omterm{def:b1:lewis-resonance-vsepr:lewis}{الأزواج الرابطة}
أكثر مما يدفع بعضها بعضًا.
\textbf{16.} الأوزون منحنٍ وذرته المركزية تختلف عن الذرتين
الطرفيتين (\omterm{def:b1:lewis-resonance-vsepr:formal-charge}{شحنة شكلية} موجبة في الوسط، وشحنة سالبة يتقاسمها
الطرفان): ثنائي قطب صغير على امتداد المنصِّف. و\ce{N2O} خطي لكن
طرفيه ذرتان مختلفتان: ثنائي قطب صغير. و\ce{N2} متناظر: لا شيء.
\textbf{17.} \ce{CO2} خطي (\ce{AX2}): تتلاشى ثنائيات أقطاب روابطه.
و\ce{SO2} منحنٍ (\ce{AX2E}): فتتجمع.
\textbf{18.} \ce{AX3E}: يضيّق \omterm{def:b1:lewis-resonance-vsepr:lewis}{الزوج الحر} الزوايا إلى ما دون
$109.5^\circ$؛ والمقيسة $106.7^\circ$.
\textbf{19.} \ce{AX2E2}: يضيّقها \omterm{def:b1:lewis-resonance-vsepr:lewis}{الزوجان الحران} أكثر؛ $104.5^\circ$.
\textbf{20.} $\mu = 2\mu_{OH}\cos(\theta/2)$
(\cref{prop:b1:lewis-resonance-vsepr:dipole-sum}).
\textbf{21.} $\mu_{OH} = 1.857/(2\cos 52.24^\circ) = 1.857/1.2246 =
\qty{1.52}{\debye}$.
\textbf{22.} $1.516 \times 3.336 \times 10^{-30} = \qty{5.06e-30}{C.m}$.
\textbf{23.} $e\,d = 1.602 \times 10^{-19} \times 95.8 \times 10^{-12} =
\qty{1.535e-29}{C.m}$، أي \qty{4.60}{\debye}.
\textbf{24.} $\delta = 5.06 \times 10^{-30}/1.535 \times 10^{-29} =
\textbf{0.33}$: تحمل الرابطة \ce{O-H} في الماء ثلث الشحنة التي
تحملها رابطة أيونية تمامًا.
\end{solution}
